% SPDX-License-Identifier: GPL-3.0-or-later
% Copyright (C) 2026 Alessandro Greco
\chapter{Onion-сервисы v3}\label{cap:14-onion-services}
\aggiornatoal{23 сентября 2026 г.}{коммит torspec \texttt{8e1f245}; tor 0.4.9.13; rend-spec v3}

\begin{obiettivi}
К концу этой главы вы сможете:
\begin{itemize}[nosep]
  \item объяснять, что предоставляет onion-сервис и чем он отличается от сайта, доступ к которому осуществляется через exit-реле;
  \item получать открытый ключ сервиса из его адреса \texttt{.onion} и выполнять обратное преобразование;
  \item описывать ослеплённые ключи, временные периоды и способ, которым клиент определяет ответственные HSDir;
  \item разбирать многоуровневую структуру дескриптора;
  \item пошагово прослеживать протокол введения и встречи (introduction и rendezvous);
  \item загружать, проверять и расшифровывать дескриптор onion-сервиса в лабораторной сети (Лабораторная работа~8).
\end{itemize}
\end{obiettivi}

До сих пор мы использовали Tor, чтобы скрывать \emph{того, кто} подключается к серверу. Сам сервер при этом оставался публичным: его IP-адрес был известен, а exit-реле подключалось к нему напрямую через публичную сеть и знало адрес назначения. Содержимое передавалось открытым текстом только в том случае, если приложение не использовало сквозное шифрование, например HTTPS. \textbf{Onion-сервисы} меняют эту схему. Они позволяют предоставлять сервис, например веб-сайт, не раскрывая его сетевое местоположение. Клиент и сервис встречаются на промежуточной точке, каждый через собственную Tor-цепь, и ни одна сторона не узнаёт адрес другой.

% ---------------------------------------------------------------------------
\section{Что предоставляет onion-сервис}\label{sec:os-proprieta}

\begin{description}
  \item[Анонимность сервиса] IP-адрес сервера никогда не раскрывается: сервис взаимодействует с сетью только через Tor-цепи, которые сам инициирует.
  \item[Самоподтверждаемые адреса] Адрес \texttt{.onion} \emph{кодирует} открытый ключ сервиса вместе с контрольной суммой и байтом версии. Любой, кто подключается по этому адресу, криптографически убеждается, что взаимодействует с владельцем соответствующего закрытого ключа, без участия центров сертификации.
  \item[Сквозное шифрование] Трафик между клиентом и сервисом шифруется ключами, известными только этим двум сторонам; exit-реле, способного наблюдать незашифрованный трафик, нет.
  \item[Доступность] Сервису не требуется принимать входящие соединения. Он может работать за NAT или межсетевым экраном, разрешающим только исходящие соединения.
\end{description}

Цена этого подхода — длина маршрута. В базовой модели клиент достигает точки встречи по трёхузловой цепи, а сервис приходит к ней по другой трёхузловой цепи: всего шесть реле вместо трёх при подключении к обычному веб-сайту. При включённом по умолчанию vanguards-lite (\cref{sec:vanguards}) маршрут перестаёт быть симметричным. Сервис добавляет фиксированное реле второго уровня, поэтому его цепь на стороне rendezvous содержит четыре перехода (\file{S\,$\to$\,G\,$\to$\,L2\,$\to$\,M\,$\to$\,Rend}), тогда как у клиента их три.

\begin{storia}[От v2 к v3]
\emph{Hidden services} существуют с 2004 года. Первая широко развёрнутая версия, v2, использовала 16-символьные адреса, полученные из 1024-битного RSA-ключа путём усечения SHA-1-дайджеста до 80 бит. Кроме того, она позволяла HSDir собирать адреса сервисов, дескрипторы которых они хранили. Версия 3 (\proposta{224}), доступная начиная с tor 0.3.2.1-alpha (2017), ввела ключи Ed25519, 56-символьные адреса, содержащие весь открытый ключ, зашифрованные дескрипторы и ослеплённые ключи. Сервисы v2 были полностью отключены в октябре 2021 года в tor 0.4.5.11 и 0.4.6.8. Изменилось и название: сегодня используется термин \emph{onion-сервис}.
\end{storia}

% ---------------------------------------------------------------------------
\section{Адрес}\label{sec:os-indirizzo}

Адрес v3 представляет собой base32-кодирование открытого идентификационного ключа сервиса, короткой контрольной суммы и номера версии~\spec{rend-spec}{Encoding onion addresses}:
\[
\texttt{address} = \mathrm{base32}(\mathit{PUBKEY} \,\|\, \mathit{CHECKSUM} \,\|\, \mathit{VERSION}) \,\|\, \texttt{".onion"}
\]
Здесь $\mathit{PUBKEY}$ — 32-байтовый ключ Ed25519, $\mathit{VERSION}$ — байт \texttt{0x03}, а $\mathit{CHECKSUM}$ состоит из первых 2 байт значения $\mathrm{SHA3\text{-}256}(\texttt{".onion checksum"} \,\|\, \mathit{PUBKEY} \,\|\, \mathit{VERSION})$. Получившиеся 35 байт кодируются ровно в 56 символов base32. Контрольная сумма не защищает от атакующего, поскольку он может вычислить её так же, как и любой другой, но позволяет обнаруживать ошибки ввода.

В одном из запусков лабораторной сети onion-сервис имел следующий адрес. Это адрес примера, сохранённого в \file{labs/ru/код/torlab/test/data/} и используемого лабораторными тестами. Поэтому все выводы в этом разделе относятся к этому примеру и могут быть воспроизведены без доступа к сети. Когда читатель запускает собственную сеть, она генерирует новые ключи и, следовательно, другой адрес (\cref{sec:rp-file}).
\begin{lstlisting}[numbers=none]
wdjtp7qktdjkkgjjbbx2axvg2gri5p4tuhlc4bq23n7o3uwu4rveppqd.onion
\end{lstlisting}
Функция \file{key\_from\_address} из Лабораторной работы~8 извлекает из этого адреса открытый ключ и проверяет его контрольную сумму.

\begin{attenzione}[Адрес нельзя надёжно узнавать по памяти]
Случайный 56-символьный адрес трудно запомнить и непросто проверить беглым взглядом. Пользователи копируют его из источника, которому доверяют, поэтому адрес, отличающийся всего несколькими символами, может остаться незамеченным. Это основной вектор \emph{фишинга} против onion-сервисов. \emph{Адреса с заданным префиксом (vanity addresses)}, создаваемые перебором миллионов ключей до тех пор, пока начальные символы адреса не образуют слово, помогают лишь частично: атакующий тоже может сгенерировать адрес с тем же префиксом. Мы вернёмся к этой проблеме в \cref{cap:21-onion-pratica}.
\end{attenzione}

% ---------------------------------------------------------------------------
\section{Ослеплённые ключи}\label{sec:os-accecate}

\subsection{Проблема}

Сервис должен где-то публиковать информацию о том, как к нему подключиться: \emph{дескриптор}. Если бы дескриптор публиковался непосредственно под идентификационным ключом сервиса, любой оператор HSDir видел бы адреса всех сервисов, чьи дескрипторы он хранит. Он мог бы перечислять такие сервисы или связывать дескрипторы одного и того же сервиса во времени. Именно так было в v2.

\subsection{Решение}

В v3 сервис подписывает дескрипторы с помощью \textbf{ослеплённого ключа}, отдельного для каждого \emph{временного периода}~\spec{rend-spec}{Appendix A: Signature scheme with key blinding}. Имея идентификационный ключ $A = aB$ (где $B$ — базовая точка Ed25519), сервис вычисляет множитель
\begin{gather*}
h = \mathrm{SHA3\text{-}256}(\texttt{"Derive temporary signing key\textbackslash 0"} \,\|\, A \,\|\, B \,\|\, N), \\
N = \texttt{"key-blind"} \,\|\, \mathit{period} \,\|\, \mathit{period\_length}
\end{gather*}
(при этом некоторые биты принудительно устанавливаются так же, как для ключей Ed25519), после чего ослеплённый ключ определяется как
\[
A' = h \cdot A.
\]
У этой конструкции есть важные свойства:
\begin{itemize}
  \item любой, кто знает адрес и, следовательно, $A$, может вычислить $A'$ для каждого периода и понять, где искать дескриптор;
  \item тот, кто видит только $A'$, например HSDir, не может восстановить $A$ или определить, что два ослеплённых ключа из разных периодов принадлежат одному сервису;
  \item сервис может подписывать данные от имени $A'$ с помощью закрытого ключа $a' = h \cdot a$, и такая подпись проверяется относительно $A'$ как обычная подпись Ed25519.
\end{itemize}

Необычная деталь хорошо показывает, насколько близко спецификация следует реализации: константная строка содержит завершающий нулевой байт, потому что tor вычисляет её длину через \texttt{sizeof}, а не \texttt{strlen}. Комментарий в исходном коде (\file{src/feature/hs/hs\_common.c}) прямо признаёт, что задокументировать этот \enquote{сюрприз} в спецификации оказалось предпочтительнее, чем менять код.

\subsection{Креденциалы и ключи, хранимые вне сети}

\begin{sloppypar}
Из идентификационного ключа выводится \emph{креденциал (credential)}: $\mathrm{SHA3\text{-}256}(\texttt{"credential"} \,\|\allowbreak\, A)$. Для каждого периода из него, в свою очередь, получается \emph{производный креденциал (subcredential)}: $\mathrm{SHA3\text{-}256}(\texttt{"subcredential"} \,\|\allowbreak\, \mathit{credential} \,\|\allowbreak\, A')$. Производный креденциал используется при шифровании дескриптора и в протоколе введения: он подтверждает знание адреса, которым HSDir не обладает.
\end{sloppypar}

Иерархия ключей также спроектирована так, чтобы идентификационный ключ мог оставаться \emph{офлайн}. Ослеплённый ключ каждого периода сертифицирует краткоживущий ключ подписи дескриптора (сертификат Ed25519 типа 8, \cref{tab:cert-types}). Теоретически сервер может заранее получать ключи подписи и сертификаты для будущих периодов, вообще не располагая ключом, соответствующим его адресу. Не все реализации поддерживают такой режим.

% ---------------------------------------------------------------------------
\section{Определение местоположения дескриптора}\label{sec:os-hsdir}

\subsection{Временные периоды}

Время делится на периоды длиной в один день в публичной сети (параметр консенсуса \texttt{hsdir-interval})~\spec{rend-spec}{Dividing time into periods}. Периоды начинаются не в полночь, а в 12:00 UTC. Они смещены на 12 часов относительно протокола совместной генерации случайных значений авторитетными серверами каталогов, который каждый день в полночь публикует в консенсусе новое \textbf{совместное случайное значение} (SRV). Номер периода вычисляется из количества минут, прошедших с 1 января 1970 года, за вычетом смещения, делённого на длину интервала. Спецификация приводит пример: в 11:15:01 UTC 13 апреля 2016 года номер временного периода равен 16903; этот период начался в полдень предыдущего дня. Лабораторная работа~8 проверяет именно этот пример.

В тестовых сетях, где голосование проводится каждые 20 секунд, ждать целый день было бы неудобно. Поэтому tor сокращает период до одного полного раунда протокола SRV, то есть до 24 голосований (функция \file{get\_time\_period\_length}). В лабораторной сети один период длится 8 минут и смещён на 4 минуты.

\subsection{Кольцо HSDir}

Какой HSDir хранит дескриптор? Каждое реле с флагом \texttt{HSDir} занимает позицию в 256-битном \emph{кольце}~\spec{rend-spec}{Where to publish a hidden service descriptor}:
\[
\mathit{relay\_index} = \mathrm{SHA3\text{-}256}(\texttt{"node-idx"} \,\|\, \textit{Ed25519 identity} \,\|\, \mathit{SRV} \,\|\, \mathit{period} \,\|\, \mathit{period\_length})
\]
Для каждой из двух \emph{реплик} сервис вычисляет индекс:
\[
\mathit{service\_index}(r) = \mathrm{SHA3\text{-}256}(\texttt{"store-at-idx"} \,\|\, A' \,\|\, r \,\|\, \mathit{period\_length} \,\|\, \mathit{period}).
\]
Сервис загружает свой дескриптор на первые четыре HSDir, следующие за каждым индексом по кольцу (параметр \file{hsdir\_spread\_store}); клиент случайным образом выбирает один из первых трёх (\file{hsdir\_spread\_fetch}). Включение SRV не позволяет атакующему задолго заранее предсказать, какие реле будут отвечать за конкретный сервис, и заранее разместить вредоносные реле в нужных позициях для цензуры сервиса или наблюдения за запросами дескриптора. Сервис одновременно публикует два дескриптора для двух последовательных периодов, чтобы оставаться доступным для клиентов, чьи консенсусы немного опережают или отстают. Какие именно два периода используются, зависит от точки в цикле. До или после ротации SRV это либо предыдущий и текущий периоды, либо текущий и следующий~\spec{rend-spec}{Generating and publishing hidden service descriptors}.

\begin{approfondimento}[Какой SRV следует использовать?]
Поскольку периоды и SRV смещены друг относительно друга на половину суток, клиент должен аккуратно выбирать SRV. В половине суток между началом периода и следующим SRV используется текущий SRV; в другой половине — предыдущий. Если в консенсусе отсутствует нужный SRV, например сразу после запуска сети, tor вычисляет запасное значение — \emph{резервное значение SRV (disaster SRV)}: $\mathrm{SHA3\text{-}256}(\texttt{"shared-random-disaster"} \,\|\, \mathit{period\_length} \,\|\, \mathit{period})$. Выбор реализован в \file{hs\_get\_current\_srv} и \file{hs\_get\_previous\_srv} (\file{src/feature/hs/hs\_common.c}). Если значение отсутствует в консенсусе, tor без дополнительного предупреждения использует резервное значение SRV. Это происходит в только что созданной лабораторной сети и длится недолго: в журнале авторитетного сервера нашей сети первый SRV появляется через 115 секунд после запуска; до этого в консенсусе нет строки \texttt{shared-rand-current-value}. Чтобы увидеть использование резервного значения, исследуйте первые несколько минут после создания сети. После некоторого времени работы сети, как в \cref{sec:lab-onion}, используются настоящие SRV.
\end{approfondimento}

\subsection{Только анонимные запросы}

Клиент загружает дескриптор HTTP-запросом к \nolinkurl{/tor/hs/3/}, за которым следует ослеплённый ключ в base64 без завершающих символов \texttt{=}, через поток \file{BEGIN\_DIR} к HSDir. Есть одно условие: запрос должен поступать по \textbf{анонимной} цепи. HSDir отклоняет запрос, пришедший по цепи, открытой клиентом напрямую, то есть по цепи из одного перехода. Мы обнаружили это в лаборатории:
\begin{lstlisting}[numbers=none]
HTTP/1.0 503 Rejecting single hop HS v3 descriptor request
\end{lstlisting}
Проверка находится в \file{connection\_dir\_is\_anonymous} (\nolinkurl{src/feature/dircommon/directory.c}). Она не позволяет HSDir связать запрос дескриптора с IP-адресом запрашивающего клиента.

% ---------------------------------------------------------------------------
\section{Дескриптор}\label{sec:os-descrittore}

Дескриптор состоит из трёх уровней~\spec{rend-spec}{Generating and publishing hidden service descriptors}, показанных на \cref{fig:os-descrittore}.

\begin{figure}[htbp]
\centering
% SPDX-License-Identifier: GPL-3.0-or-later
% Copyright (C) 2026 Alessandro Greco
%% Русское альтернативное описание EPUB.
\altfig{Трёхуровневый дескриптор onion v3: внешняя подписанная оболочка видна HSDir и содержит версию, срок действия, сертификат типа 8, ревизию, superencrypted и подпись Ed25519. Первый слой шифруется с помощью ослеплённого открытого ключа и производного креденциала, содержит desc-auth-type, desc-auth-ephemeral-key, auth-client и encrypted. Даже без ограниченного обнаружения существуют фиктивные записи auth-client. Внутренний слой хранит параметры рукопожатий, управления потоком и ключи точек введения. При ограниченном обнаружении закрытый ключ X25519 разрешённого клиента извлекает cookie дескриптора из auth-client, которым расшифровывают внутренний слой. Числа приведены как примеры.}
% Слои дескриптора
\begin{tikzpicture}[font=\sffamily\scriptsize,
  strato/.style={draw,rounded corners=3pt,align=left,inner sep=5pt,text width=#1},
  lab/.style={font=\sffamily\scriptsize\bfseries,anchor=north west}]
  \node[strato=10.5cm,fill=torgrigio!8,draw=torgrigio,minimum height=6.3cm] (e) at (0,0) {};
  \node[lab] at ([shift={(0.15,-0.1)}]e.north west) {Внешний слой (открытый, подписан)};
  \node[anchor=north west,align=left,font=\ttfamily\scriptsize] at ([shift={(0.3,-0.5)}]e.north west)
    {hs-descriptor 3\\descriptor-lifetime 180\\descriptor-signing-key-cert \normalfont\sffamily(тип 8, от ослеплённого ключа)\\\ttfamily revision-counter 662427\\superencrypted};
  \node[strato=9.6cm,fill=torviola!10,draw=torviola,minimum height=3.35cm] (s) at (0.35,-0.85) {};
  \node[lab] at ([shift={(0.15,-0.1)}]s.north west) {Первый слой: ослеплённый ключ + подкреденциал};
  \node[anchor=north west,align=left,font=\ttfamily\scriptsize] at ([shift={(0.3,-0.5)}]s.north west)
    {desc-auth-type, desc-auth-ephemeral-key, auth-client \ldots\\encrypted};
  \node[strato=8.7cm,fill=torverde!15,draw=torverde!60!black,minimum height=1.75cm] (i) at (0.7,-1.35) {};
  \node[lab] at ([shift={(0.15,-0.1)}]i.north west) {Второй слой (cookie клиента при ограничении)};
  \node[anchor=north west,align=left,font=\ttfamily\scriptsize] at ([shift={(0.3,-0.5)}]i.north west)
    {create2-formats 2 \quad flow-control 1-2 31\\introduction-point \ldots\ onion-key ntor \ldots\ auth-key \ldots\ enc-key \ldots};
  \node[anchor=south west,font=\ttfamily\scriptsize] at ([shift={(0.3,0.15)}]e.south west) {signature \normalfont\sffamily(Ed25519, префикс \texttt{"Tor onion service descriptor sig v3"})};
\end{tikzpicture}

\caption{Структура дескриптора onion-сервиса v3. Подписанная внешняя оболочка читается HSDir; первый зашифрованный слой раскрывается при знании идентичности сервиса. При ограниченном обнаружении разрешённый клиент использует ключ X25519 для извлечения cookie дескриптора, необходимого для расшифрования второго слоя.}
\label{fig:os-descrittore}
\end{figure}

\paragraph{Внешний уровень} передаётся открытым текстом и подписывается. Он содержит версию, срок действия (\texttt{descriptor-lifetime}, 180 минут в лабораторной сети), сертификат ключа подписи, выданный ослеплённым ключом, счётчик ревизий, зашифрованный блок \texttt{superencrypted} и подпись. Подпись — Ed25519, созданная ключом подписи дескриптора по тексту с префиксом \texttt{"Tor onion service descriptor sig v3"}. Префикс не позволяет подписи, созданной для дескриптора, оказаться допустимой для другого типа данных. Счётчик ревизий должен увеличиваться с каждой новой версией, чтобы HSDir не принимал дескрипторы старее уже сохранённых.

\paragraph{Первый зашифрованный уровень} шифруется ключом, выведенным из ослеплённого ключа и производного креденциала. HSDir, не знающий адреса, не может прочитать этот уровень. В нём находятся поля \emph{ограниченного обнаружения (restricted discovery)} (\cref{sec:os-ristretta}) и второй уровень. Чтобы не раскрывать сам факт включения ограниченного обнаружения, эти поля присутствуют всегда и заполняются случайными данными, когда они не нужны.

\paragraph{Второй уровень} содержит сведения, необходимые для подключения: поддерживаемые форматы рукопожатия, параметры управления потоком, возможные параметры защиты от отказа в обслуживании (DoS) и, для каждой \textbf{точки введения} --- спецификаторы соединения для доступа к ней (в base64), ключ ntor, сертифицированный ключ аутентификации и ключ шифрования. По умолчанию сервис использует три точки введения.

Оба уровня используют одну и ту же схему. Случайная 16-байтовая соль, секретные данные и константа, специфичная для уровня, подаются на SHAKE-256, чтобы вывести ключ AES-256, вектор инициализации и ключ аутентификации. Текст шифруется в режиме счётчика и аутентифицируется с помощью SHA3-256. Открытый текст первого уровня дополняется до длины, кратной 10\,000 байт, а строки \texttt{auth-client} — до количества, кратного 16, чтобы размер дескриптора раскрывал как можно меньше информации о числе точек введения и авторизованных клиентов.

% ---------------------------------------------------------------------------
\section{Введение и встреча (rendezvous)}\label{sec:os-rendezvous}

Получив дескриптор, клиент может подключиться. Протокол, показанный на \cref{fig:os-rendezvous}, использует три \enquote{специальные} роли: точку введения, выбранную сервисом, точку встречи, выбранную клиентом, и серверы HSDir.

\begin{figure}[htbp]
\centering
% SPDX-License-Identifier: GPL-3.0-or-later
% Copyright (C) 2026 Alessandro Greco
%% Русское альтернативное описание EPUB.
\altfig{Протокол встречи onion-сервиса: сервис строит многоузловую цепь к точке введения и передаёт ESTABLISH\_INTRO (1), затем публикует дескриптор в HSDir (2). Клиент загружает его, строит отдельную цепь к узлу встречи и отправляет ESTABLISH\_RENDEZVOUS с cookie (3); через другой маршрут к точке введения он посылает INTRODUCE1 (4). Точка введения пересылает сообщение сервису как INTRODUCE2 по ранее созданной сервисом цепи, не читая зашифрованное содержимое. Сервис отправляет RENDEZVOUS1 к узлу встречи (5), который связывает две цепи и пересылает клиенту RENDEZVOUS2 (6). Все линии изображают логические многоузловые пути, не прямое соединение клиента и сервиса.}
% Протокол встречи
\begin{tikzpicture}[font=\sffamily\small,>=Stealth,
  circ/.style={-{Stealth[length=2.5mm]},thick,decorate,decoration={snake,amplitude=.5mm,segment length=4mm,post length=3mm}},
  msg/.style={font=\ttfamily\scriptsize,fill=white,inner sep=1.5pt,align=center}]
  \node[client] (c) at (0,0) {Клиент};
  \node[client,fill=torverde!15,draw=torverde!60!black] (s) at (11,0) {Онион-сервис};
  \node[relay,fill=torgrigio!10,draw=torgrigio,font=\sffamily\scriptsize] (h) at (1.2,3.8) {HSDir};
  \node[relay,fill=torviola!15,draw=torviola,align=center,font=\sffamily\scriptsize] (ip) at (8.2,3.6) {точка\\введения};
  \node[relay,fill=torverde!25,align=center,font=\sffamily\scriptsize] (rp) at (5.5,-3) {узел\\встречи};
  % Сервис
  \draw[circ,torverde!60!black] (s) -- node[msg,pos=0.35,right=12pt]{(1) ESTABLISH\_INTRO} (ip);
  \draw[circ,torverde!60!black] (s.north) to[out=80,in=20,looseness=1.1] node[msg,pos=0.45,above=2pt]{(2) публикует дескриптор} (h.north east);
  % Клиент
  \draw[circ,torviola] (c) -- node[msg,pos=0.5,left=4pt]{загружает\\дескриптор} (h);
  \draw[circ,torviola] (c) -- node[msg,pos=0.45,left=6pt]{(3) ESTABLISH\_RENDEZVOUS\\(cookie)} (rp);
  \draw[circ,torviola] (c.north east) to[out=15,in=210] node[msg,pos=0.5,below=3pt]{(4) INTRODUCE1} (ip.south west);
  \draw[->,thick,torviola,dashed] (ip.south east) to[out=-30,in=100] node[msg,pos=0.45,left=3pt]{INTRODUCE2} (s.north);
  \draw[circ,torverde!60!black] (s) -- node[msg,pos=0.45,right=6pt]{(5) RENDEZVOUS1\\(cookie)} (rp);
  \node[font=\sffamily\scriptsize,align=center,below=3pt of rp] {(6) связывает две цепи\\и передаёт RENDEZVOUS2 клиенту};
\end{tikzpicture}

\caption{Протокол встречи (rendezvous). Каждая стрелка представляет многоузловую цепь; прямого соединения между клиентом, сервисом и промежуточными точками нет. Числа показывают порядок шагов, описанных в тексте.}
\label{fig:os-rendezvous}
\end{figure}

\begin{enumerate}
  \item \textbf{Сервис подготавливает точки введения.} Он выбирает несколько реле, строит цепь к каждому и отправляет \file{ESTABLISH\_INTRO} (команда реле 32), содержащую отдельный ключ аутентификации и подпись для каждой точки. Реле отвечает \file{INTRO\_ESTABLISHED} (38); после этого оно пересылает сервису запросы, содержащие соответствующий ключ.
  \item \textbf{Сервис публикует дескриптор} на ответственных HSDir, как описано выше.
  \item \textbf{Клиент подготавливает точку встречи.} Он случайным образом выбирает реле в качестве точки встречи, строит к нему цепь и отправляет \file{ESTABLISH\_RENDEZVOUS} (33) со случайным 20-байтовым \emph{идентификатором (cookie)}. Реле отвечает \file{RENDEZVOUS\_ESTABLISHED} (39).
  \item \textbf{Клиент представляет себя сервису.} Он строит цепь к точке введения и отправляет \texttt{INTRODUCE1} (34). Открытая часть указывает ключ аутентификации точки введения. Зашифрованная часть, доступная только сервису, содержит идентификатор cookie, спецификаторы соединения и ключ ntor точки встречи, а также первую часть рукопожатия. Точка введения отправляет клиенту \file{INTRODUCE\_ACK} (40) и пересылает сообщение сервису как \texttt{INTRODUCE2} (35).
  \item \textbf{Сервис достигает точки встречи.} Он расшифровывает \texttt{INTRODUCE2}, проверяет, не дублирует ли сообщение уже обработанный запрос (\emph{кеш защиты от повторов (replay cache)}), строит цепь к указанной точке встречи и отправляет \texttt{RENDEZVOUS1} (36) с идентификатором cookie и второй частью рукопожатия.
  \item \textbf{Точка встречи соединяет цепи.} Она находит клиентскую цепь по тому же идентификатору cookie и пересылает вторую часть рукопожатия в сообщении \texttt{RENDEZVOUS2} (37), после чего начинает копировать ячейки между двумя цепями.
\end{enumerate}

Рукопожатие между клиентом и сервисом, называемое \emph{hs-ntor}, является вариантом ntor (\cref{sec:ntor}), в который включены производный креденциал и ключ аутентификации точки введения. Только сторона, знающая адрес, может сформировать \texttt{INTRODUCE1}, которое сервис примет. Выведенные ключи добавляют ещё один слой шифрования, подобно \emph{виртуальному переходу}, добавленному к клиентской цепи, и используют SHA3-256 для накопительного дайджеста вместо SHA-1 (\cref{sec:riconoscere}). Ни точка введения, ни точка встречи не могут читать трафик.

Через объединённую цепь клиент открывает потоки командой \texttt{BEGIN}, как при работе через exit-реле (\cref{cap:11-stream}). Сервис перенаправляет эти потоки на локальные порты, настроенные директивой \opt{HiddenServicePort}.

\begin{approfondimento}[Зачем нужно столько реле?]
Точка введения знает сервис только как \enquote{ключ аутентификации некоторой цепи}; она не знает, какому адресу \texttt{.onion} служит. Точка встречи видит случайный идентификатор cookie и две цепи; она не знает, какой сервис соединяет с каким клиентом. Каждая роль видит лишь фрагмент картины, и ни одна не знает IP-адрес клиента или сервиса, которые остаются скрыты за соответствующими входными узлами (guards).
\end{approfondimento}

% ---------------------------------------------------------------------------
\section{Ограниченное обнаружение (restricted discovery)}\label{sec:os-ristretta}

Сервис может быть предназначен только для определённых клиентов, например если это административная панель. При использовании \textbf{ограниченного обнаружения (restricted discovery)} (ранее \emph{авторизации клиентов}) каждый авторизованный клиент имеет пару ключей X25519, а сервис знает её открытую часть. В строках \texttt{auth-client} первого уровня сервис отдельно шифрует для каждого клиента ключ второго уровня дескриптора. Тот, кто знает адрес, но не располагает авторизованным ключом, не может прочитать список точек введения и, следовательно, даже попытаться подключиться. В tor это настраивается добавлением открытых ключей клиентов в каталог \file{authorized\_clients} сервиса.

% ---------------------------------------------------------------------------
\section{Защитные механизмы}\label{sec:os-difese}

Onion-сервисы являются естественными целями атак: сервис не может просто блокировать IP-адреса атакующих клиентов, поскольку не видит их. Основные защитные механизмы, подробнее рассматриваемые в других разделах, таковы:
\begin{itemize}
  \item \textbf{vanguards}, защищающие guard-реле сервиса от атак, основанных на цепях, создаваемых по запросу (\cref{sec:vanguards});
  \item \textbf{кеш защиты от повторов (replay cache)}, не позволяющий атакующему заставлять сервис повторно выполнять одну и ту же работу;
  \item параметры защиты от DoS, которые сервис может задавать на точках введения через расширение \file{ESTABLISH\_INTRO}, ограничивая количество сообщений \texttt{INTRODUCE2} в секунду;
  \item \textbf{доказательство выполнения работы (proof of work)}, при котором атакуемый сервис требует от клиентов решить небольшую вычислительную задачу перед обслуживанием (\cref{cap:15-dos}).
\end{itemize}

% ---------------------------------------------------------------------------
\section{Лабораторная работа: поиск и чтение дескриптора}\label{sec:lab-onion}

\begin{laboratorio}[8]{Дескриптор onion-сервиса лабораторной сети}
\textbf{Цель.} Реализовать \file{labs/ru/код/torlab/onion.py}: обработку адресов, ослеплённого ключа, кольца HSDir, загрузку, проверку и расшифрование дескриптора.\\
\textbf{Предварительные условия.} Лабораторные работы~1--6; лабораторная сеть должна быть запущена.\\
\textbf{Оценочное время.} 4 часа.
\end{laboratorio}

\paragraph{Шаг 1: ослеплённый ключ.} Распространённые криптографические библиотеки поддерживают подписи Ed25519, но не умножение точки на скаляр, необходимое для вычисления $A' = h \cdot A$. Поэтому модуль содержит небольшую реализацию арифметики Ed25519 — около тридцати строк, основанных на формулах RFC~8032. Она достаточна для этой задачи, но непригодна для общего использования: работает медленно и не защищает от атак по времени выполнения.
\begin{lstlisting}[language=Python]
def blinding_factor(A, period, duration_minutes):
    n = b"key-blind" + struct.pack(">QQ", period, duration_minutes)
    h = bytearray(hashlib.sha3_256(
        b"Derive temporary signing key\x00" + A + BASE_STR + n).digest())
    h[0] &= 248; h[31] &= 63; h[31] |= 64
    return int.from_bytes(h, "little")

def blinded_key(A, period, duration_minutes):
    return _encode_point(_multiply(blinding_factor(A, period, duration_minutes),
                                   _decode_point(A)))
\end{lstlisting}

\paragraph{Шаг 2: период и SRV.} Функция \texttt{time\_period\_info} воспроизводит вычисление tor на основе метки времени \texttt{valid-after} и интервала голосования из консенсуса, определяя, нужен текущий или предыдущий SRV. Если в консенсусе нет требуемого значения, используется резервное значение SRV.

\paragraph{Шаг 3: трассировка ошибки.} При первом запуске наш код не нашёл дескриптор ни на одном из восьми HSDir сети. Поскольку мы проверили их все, ошибка выбора HSDir не могла объяснить проблему. Оставались две возможности: неверный ослеплённый ключ или неверный запрос. Чтобы разделить их, мы заставили настоящий клиент сети загрузить дескриптор через управляющий порт (\texttt{HSFETCH}, затем \texttt{GETINFO hs/client/desc/id/...}). Из полученного сертификата мы прочитали ослеплённый ключ, реально использованный сервисом: он совпал с нашим. Следовательно, проблема была в запросе, и исходный код tor показал причину: HSDir отклоняют запросы по односкачковым цепям. После построения трёхузловой цепи к HSDir мы получили дескриптор. Это универсальный метод: когда эксперимент не работает, ищите промежуточную точку, в которой ваш результат можно сравнить с эталонной реализацией.

\paragraph{Шаг 4: проверка и расшифрование.} В примере \texttt{path} обозначает список из трёх объектов \texttt{RelayInfo}, полученных из лабораторной сети, с HSDir последним; \texttt{A} --- ключ идентичности, \texttt{Ap} --- ослеплённый ключ для этого периода. Создание объектов здесь не показано.
\begin{lstlisting}[language=Python]
d = download_descriptor(Ap, path)
desc = parse_descriptor(d, A, Ap)
\end{lstlisting}
\file{parse\_descriptor} проверяет, что сертификат ключа подписи был выдан ожидаемым ослеплённым ключом, проверяет подпись дескриптора и расшифровывает оба уровня. Сокращённый результат предыдущей лабораторной сессии:
\begin{lstlisting}[numbers=none]
срок 180 ревизия 662427 точек введения 3
create2-formats 2
flow-control 1-2 31
introduction-point AwAGfwAAARPyAhQXE5GoH+wp3MoxW9M1QV6XSxb0fgM...
onion-key ntor yRURxDb2eeEmF/p5M/bmHA+giwVyxjApKYyNY1KkCDo=
auth-key
-----BEGIN ED25519 CERT-----
...
\end{lstlisting}
Строка \texttt{flow-control 1-2 31} сообщает клиентам, что сервис поддерживает управление перегрузкой с тем же значением \file{sendme\_inc}, равным 31, которое встречалось в \cref{cap:11-stream}.

\paragraph{Шаг 5: проверка согласованности.} В лабораторной сети восемь HSDir, а сервис загружает дескриптор на $2 \times 4 = 8$ из них, то есть на все. Поэтому само обнаружение дескриптора ещё не доказывает правильность вычисления кольца. Однако журнал клиента (\file{info.log}) фиксирует HSDir, к которому он обращался: \texttt{test002a} в 16:07. Это одно из трёх реле, возвращённых нашим кодом для данной реплики в тот момент. Одно наблюдение не является доказательством, но согласуется с расчётом. В упражнении~\ref{es:14-anello} предлагается более строгий тест.

% ---------------------------------------------------------------------------
\section*{Итоги}
\begin{itemize}
  \item Onion-сервис предоставляет сервис, не раскрывая его IP-адрес; адрес \texttt{.onion} содержит и аутентифицирует открытый ключ сервиса.
  \item Дескрипторы подписываются ослеплёнными ключами, специфичными для периода; связать такой ключ с сервисом может только сторона, знающая адрес.
  \item Позиции дескриптора в кольце HSDir зависят от ослеплённого ключа, периода и совместного случайного значения; клиенты загружают дескрипторы через анонимные запросы.
  \item Дескриптор содержит подписанный внешний уровень и два зашифрованных уровня; второй перечисляет точки введения.
  \item Клиент и сервис встречаются в точке встречи, выбранной клиентом, после запроса через точку введения. Рукопожатие hs-ntor добавляет слой сквозного шифрования.
  \item Ограниченное обнаружение, vanguards, кеш повторных запросов, ограничение частоты запросов и доказательство выполнения работы защищают сервисы.
\end{itemize}

% ---------------------------------------------------------------------------
\section*{Упражнения}

\begin{esercizio}[1]\label{es:14-indirizzo}
Сколько байт кодирует адрес v3 и почему в результате получается ровно 56 символов base32? Какова вероятность того, что случайная ошибка ввода пройдёт проверку контрольной суммы?
\end{esercizio}

\begin{esercizio}[1]\label{es:14-periodo}
Вычислите номер временного периода публичной сети для 09:00 UTC 1 января 2026 года. Когда начался этот период?
\end{esercizio}

\begin{esercizio}[2]\label{es:14-hsdir}
Атакующий хочет заблокировать доступ к onion-сервису, добившись контроля над ответственными HSDir. Что ему необходимо предсказать и насколько заранее это возможно? Какую роль играет SRV?
\end{esercizio}

\begin{esercizio}[2]\label{es:14-ruoli}
Перечислите, что точка введения, точка встречи и HSDir знают о клиенте и сервисе. Какая из этих ролей была бы наиболее интересна атакующему, пытающемуся оценить популярность сервиса?
\end{esercizio}

\begin{esercizio}[2]\label{es:14-anonimo}
Почему HSDir отклоняет запросы дескриптора, поступающие по односкачковым цепям? Что в противном случае он мог бы узнать?
\end{esercizio}

\begin{esercizio}[3]\label{es:14-anello}
Проверьте вычисление кольца более строго. Добавьте реле в лабораторную сеть (упражнение~\ref{es:19-rete}), пока в ней не станет больше восьми HSDir, а затем убедитесь, что дескриптор доступен \emph{только} на HSDir, выбранных вашим кодом.
\end{esercizio}

\begin{esercizio}[3]\label{es:14-ipt}
Декодируйте спецификаторы соединения точки введения из лабораторного дескриптора (значение base64 в строке \texttt{introduction-point}). Какому реле лабораторной сети она соответствует? Что содержит сертификат в строке \texttt{auth-key}?
\end{esercizio}

% ---------------------------------------------------------------------------
\section*{Дополнительная литература}
\begin{itemize}
  \item \titolo{Tor Rendezvous Specification - Version 3} (rend-spec)~\cite{torspec}: адреса, ослеплённые ключи, дескрипторы, introduction и rendezvous, включая тестовые векторы hs-ntor.
  \item Предложение 224: архитектура сервисов v3.
  \item Спецификация совместной генерации случайных значений (srv-spec): как авторитетные серверы каталогов формируют SRV.
\end{itemize}
